%% kapa: exam=2 round=23 year=2012 period=3 subject=compensation_law subject_ko=감정평가 및 보상법규 seq=3 number=3 points=20
%% source: Q-Net 감정평가사 기출문제 (cst00309) artlSeq=5056014 file=raw/1008/5056014/3교시(감정평가 및 보상법규).hwp
%% conversion: kordoc 4.24.51 -> tools/gen2.py (auto); review=proofread-sonnet; flags=none
\begin{longquestion}{【문제 3】}
20년 이상 감정평가업에 종사하고 있는 감정평가사 甲은 2년 전에 국토해양부장관 乙의 인가를 받아 50명 이상의 종업원을 고용하는 감정평가법인을 설립하였다. 그 후 乙은 甲이 정관을 거짓으로 작성하는 등 부정한 방법으로 감정평가법인의 설립인가를 받았다는 이유로, 「부동산 가격공시 및 감정평가에 관한 법률」제38조 제1항 제6호에 따라 설립인가를 취소하였다. 甲은 乙의 인가취소가 잘못된 사실관계에 기초한 위법한 처분이라는 이유로 취소소송을 제기하면서 집행정지신청을 하였다. 甲의 집행정지신청의 인용여부를 논하시오. (20점) \par
\end{longquestion}
